%%% FIELD psc-kill-statement
Determine the exact value of \(ID_{K_6}\) by an exact, reproducible computation of \(\operatorname{Cl}_d(K_6)\) for \(d=1,\ldots,9\).  The artifact must contain a deterministic script, a hash-addressed result, and, at every recursive round and for every structural or disturbance coordinate not yet adjoined, the exact monomial-image matrices or equivalent pivot and nullspace certificates for the intersections of \(A_d(S)\) with \(qC_d(S)\).  The least \(d\) with full closure is \(ID_{K_6}\) by \(\mathrm{thm:closure\mbox{-}characterization}\).  The proved bound \(ID_{K_6}\leq7\) requires full closure by degree seven; the degree-eight and degree-nine records are monotonicity checks and compatibility checks with the earlier degree-nine construction.  No conclusion about auxiliary-cutset certification or a degree-seven lower bound is part of this question.
%%% FIELD psc-frontier-statement
On the unchanged class \(\mathcal G_{2}^{\mathrm{rat}}(n)\), determine the full-class elbow
\[
 \sup_{n\geq1}D_2(n)<\infty\quad\text{or}\quad \sup_{n\geq1}D_2(n)=\infty.
\]
In the bounded case, determine the least uniform constant and exhibit all required exact closure witnesses.  In the unbounded case, determine the sharp growth rate and construct an explicit dimension-feasible family \(G_n\in\mathcal G_{2}^{\mathrm{rat}}(n)\) with an order-robust proper degree closure at the matching lower scale.  Every proposed lower family must first satisfy
\[
 \binom n2-|B|\geq |D|,
\]
must be proved generically birational, and must control recursive substitutions through every structural and disturbance coordinate.  The unconditional bound \(D_2(n)\leq U_2(n)\) is available but is not asserted to be sharp.  The graphs \(W_{\ell}^{(2)}\) cannot serve as the lower family because \(\mathrm{thm:windmill\mbox{-}overparameterized}\) proves that they are generically infinite-to-one.
%%% FIELD psc-stopping-statement
Let \(G\) be any acyclic mixed graph on \(n\) vertices with directed indegree at most two, and set \(B_2(n)=U_2(n)\).  One exact call to \(DBI\) at degree \(B_2(n)\), with no wall-clock timeout, terminates and returns ``yes'' with minimum-degree rational formulas if and only if \(G\in\mathcal G_{2}^{\mathrm{rat}}(n)\).  Its arithmetic complexity is
\[
 O\!\left(B_2(n)n^5\{3B_2(n)n\}^{4\alpha_{\mathrm{alg}}B_2(n)^2\log(3B_2(n)n)}\right).
\]
A ``no'' answer is accompanied by a complete exact bounded-search certificate rather than a timeout status.  This schedule is unconditional, although no sharpness claim is made for \(B_2(n)\).
%%% FIELD pdc-frontier-construction
Apply \(\mathsf{FrontierHandle}\) to the unchanged class \(\mathcal G_{2}^{\mathrm{rat}}(n)\) by starting from the proved effective bound \(D_2(n)\leq U_2(n)\) and deciding whether \(\sup_nD_2(n)\) is finite.  In the bounded branch, determine the least uniform constant by exact degree closures.  In the unbounded branch, construct a generically birational indegree-two family and prove matching upper and lower rates, first checking the dimension-feasibility inequality \(\binom n2-|B|\geq|D|\) and then exhibiting a proper \(\operatorname{Cl}_d(G)\) that remains proper after every permitted earlier adjunction of structural or disturbance coordinates.
%%% FIELD edit-frontier-construction
Apply \(\mathsf{FrontierHandle}\) to the unchanged class \(\mathcal G_{2}^{\mathrm{rat}}(n)\) by starting from the proved effective bound \(D_2(n)\leq U_2(n)\) and deciding whether \(\sup_nD_2(n)\) is finite.  In the bounded branch, determine the least uniform constant by exact degree closures.  In the unbounded branch, construct a generically birational indegree-two family and prove matching upper and lower rates, first checking the dimension-feasibility inequality \(\binom n2-|B|\geq|D|\) and then exhibiting a proper \(\operatorname{Cl}_d(G)\) that remains proper after every permitted earlier adjunction of structural or disturbance coordinates.
%%% FIELD pdc-stopping-construction
Compute \(U_2(n)=2\cdot8^{2^{n(n+1)}}\), set \(B_2(n)=U_2(n)\), run \(DBI\) once through its complete internal weighted-degree loop at cutoff \(B_2(n)\), disable finite wall-clock timeouts, and retain the exact identifying formulas or the complete exact bounded-search certificate.
%%% FIELD edit-stopping-construction
Compute \(U_2(n)=2\cdot8^{2^{n(n+1)}}\), set \(B_2(n)=U_2(n)\), run \(DBI\) once through its complete internal weighted-degree loop at cutoff \(B_2(n)\), disable finite wall-clock timeouts, and retain the exact identifying formulas or the complete exact bounded-search certificate.
%%% FIELD edit-kill-statement
Determine the exact value of \(ID_{K_6}\) by an exact, reproducible computation of \(\operatorname{Cl}_d(K_6)\) for \(d=1,\ldots,9\).  The artifact must contain a deterministic script, a hash-addressed result, and, at every recursive round and for every structural or disturbance coordinate not yet adjoined, the exact monomial-image matrices or equivalent pivot and nullspace certificates for the intersections of \(A_d(S)\) with \(qC_d(S)\).  The least \(d\) with full closure is \(ID_{K_6}\) by \(\mathrm{thm:closure\mbox{-}characterization}\).  The proved bound \(ID_{K_6}\leq7\) requires full closure by degree seven; the degree-eight and degree-nine records are monotonicity checks and compatibility checks with the earlier degree-nine construction.  No conclusion about auxiliary-cutset certification or a degree-seven lower bound is part of this question.
%%% FIELD edit-frontier-statement
On the unchanged class \(\mathcal G_{2}^{\mathrm{rat}}(n)\), determine the full-class elbow
\[
 \sup_{n\geq1}D_2(n)<\infty\quad\text{or}\quad \sup_{n\geq1}D_2(n)=\infty.
\]
In the bounded case, determine the least uniform constant and exhibit all required exact closure witnesses.  In the unbounded case, determine the sharp growth rate and construct an explicit dimension-feasible family \(G_n\in\mathcal G_{2}^{\mathrm{rat}}(n)\) with an order-robust proper degree closure at the matching lower scale.  Every proposed lower family must first satisfy
\[
 \binom n2-|B|\geq |D|,
\]
must be proved generically birational, and must control recursive substitutions through every structural and disturbance coordinate.  The unconditional bound \(D_2(n)\leq U_2(n)\) is available but is not asserted to be sharp.  The graphs \(W_{\ell}^{(2)}\) cannot serve as the lower family because \(\mathrm{thm:windmill\mbox{-}overparameterized}\) proves that they are generically infinite-to-one.
%%% FIELD edit-stopping-statement
Let \(G\) be any acyclic mixed graph on \(n\) vertices with directed indegree at most two, and set \(B_2(n)=U_2(n)\).  One exact call to \(DBI\) at degree \(B_2(n)\), with no wall-clock timeout, terminates and returns ``yes'' with minimum-degree rational formulas if and only if \(G\in\mathcal G_{2}^{\mathrm{rat}}(n)\).  Its arithmetic complexity is
\[
 O\!\left(B_2(n)n^5\{3B_2(n)n\}^{4\alpha_{\mathrm{alg}}B_2(n)^2\log(3B_2(n)n)}\right).
\]
A ``no'' answer is accompanied by a complete exact bounded-search certificate rather than a timeout status.  This schedule is unconditional, although no sharpness claim is made for \(B_2(n)\).
%%% FIELD edit-inference-statement
Fix \(n\), \(G\), \(r\), and a deterministic returned map \(F_G\), assume that \(\mathcal C_G(\delta_{\mathrm{chart}},\kappa_{\mathrm{chart}},K_{\mathrm{chart}},\nu_{\mathrm{chart}})\neq\varnothing\), and use the total pole and singular-covariance conventions in \(\mathrm{def:plug\mbox{-}in\mbox{-}estimator}\) and \(\mathrm{def:simultaneous\mbox{-}wald\mbox{-}set}\). Uniformly over \(\Sigma\in\mathcal C_G(\delta_{\mathrm{chart}},\kappa_{\mathrm{chart}},K_{\mathrm{chart}},\nu_{\mathrm{chart}})\),
\[
\sqrt N\{\widehat\theta_G-F_G(s)\}=N^{-1/2}\sum_{i=1}^N\operatorname{IF}_F(X_i)+o_{\mathbb P}(1)\rightsquigarrow\mathcal N(0,V_F),\qquad \widehat V_F\xrightarrow{\mathbb P}V_F.
\]
Moreover,
\[
\sup_{\Sigma\in\mathcal C_G}\Pr_{\Sigma}\!\left\{\min_{1\le j\le r}|Q_{F,j}(\operatorname{vech}\widehat\Sigma)|=0\right\}\to0,\qquad \sup_{\Sigma\in\mathcal C_G}\Pr_{\Sigma}\{\widehat V_F\text{ is singular}\}\to0.
\]
Consequently \(\inf_{\Sigma\in\mathcal C_G}\Pr_{\Sigma}\{F_G(s)\in\mathcal W_{G,N}(1-\gamma)\}\to1-\gamma\). The result covers every fixed finite vector of rational formulas returned by the total schedule on the Gaussian submodel; it makes no efficiency claim.
%%% FIELD cubic-statement
For every acyclic mixed graph \(G=(V_n,D,B)\), put the formal covariance coordinates \(s\) into a symmetric matrix \(\Sigma\), and interpret \(\omega_{ij}=0\) whenever \(i\neq j\) and \(i\leftrightarrow j\notin B\).  Then
\[
 I_G=\left\langle
 \bigl[(I_n-\Lambda)^{\mathsf T}\Sigma(I_n-\Lambda)\bigr]_{ij}-\omega_{ij}:1\leq i\leq j\leq n
 \right\rangle
 \subseteq\mathbb Q[\theta_G,s].
\]
In particular, \(I_G\) has a generating set of total degree at most three.
%%% FIELD cubic-proof
Let \(R=I_n-\Lambda\).  By acyclicity, after a simultaneous permutation of rows and columns, \(\Lambda\) is strictly triangular.  Hence \(\Lambda^n=0\), \(\det R=1\), and
\[
 T:=R^{-1}=I_n+\Lambda+\cdots+\Lambda^{n-1}
\]
has entries in \(\mathbb Q[\theta_G]\).  Let \(J\) denote the displayed ideal, with each structurally absent off-diagonal entry of \(\Omega\) set equal to zero as required by \(\mathrm{ass:omega\mbox{-}support}\).  In the quotient by \(J\), the defining relations are
\[
 R^{\mathsf T}\Sigma R=\Omega.
\]
Multiplication on the left by \(T^{\mathsf T}\) and on the right by \(T\) is valid over the polynomial ring and gives
\[
 \Sigma=T^{\mathsf T}\Omega T.
\]
Consequently the homomorphism
\[
 \mathbb Q[\theta_G,s]/J\longrightarrow\mathbb Q[\theta_G]
\]
which fixes \(\theta_G\) and sends \(s\) to \(\phi_G(\theta_G)\) is an isomorphism: its inverse sends every source coordinate to its residue class, and the last display shows that every covariance residue class is the prescribed polynomial in those source coordinates.  Its kernel before quotienting is therefore exactly \(J\).  By the definition \(I_G=\ker\tau_G\), this proves \(I_G=J\).

Each entry of \(R^{\mathsf T}\Sigma R-\Omega\) is a sum of terms of the forms \(\sigma_{ij}\), \(\lambda_{ab}\sigma_{ij}\), \(\lambda_{ab}\sigma_{ij}\lambda_{cd}\), and \(\omega_{ij}\).  Their total degrees are at most three, proving the final assertion.
%%% FIELD dube-statement
Let \(n\geq1\).  An ideal over a field in at most \(n(n+1)\) indeterminates, generated by polynomials of total degree at most three, has for every admissible monomial order a reduced Gr\"obner basis whose elements have total degree at most
\[
 2\left(\frac{15}{2}\right)^{2^{n(n+1)-1}}<U_2(n)=2\cdot8^{2^{n(n+1)}}.
\]
%%% FIELD upper-statement
For every \(n\geq1\) and every \(G\in\mathcal G_{2}^{\mathrm{rat}}(n)\),
\[
 ID_G\leq U_2(n)=2\cdot8^{2^{n(n+1)}}.
\]
More strongly, for every \(q\in\theta_G\) there are \(a_q,b_q\in\mathbb Q[s]\) such that \(\tau_G(a_q)\neq0\),
\[
 q a_q(s)-b_q(s)\in I_G,
 \qquad
 \deg\{q a_q-b_q\}\leq U_2(n).
\]
Thus every source coordinate has a direct certificate of the displayed degree, without recursive use of previously identified source coordinates, and \(D_2(n)\leq U_2(n)\).
%%% FIELD upper-proof
Fix \(G\in\mathcal G_{2}^{\mathrm{rat}}(n)\), and write
\[
 p=|\theta_G|,\qquad m=\frac{n(n+1)}2.
\]
Generic birationality of \(\phi_G\), which is part of \(\mathrm{def:graph\mbox{-}class}\), identifies the function field generated by the \(m\) covariance coordinates with \(\mathbb Q(\theta_G)\).  Taking transcendence degrees over \(\mathbb Q\) gives
\[
 p=\operatorname{trdeg}_{\mathbb Q}\mathbb Q(\theta_G)
   =\operatorname{trdeg}_{\mathbb Q}\mathbb Q(\phi_G(\theta_G))
   \leq m. \tag{1}
\]
Therefore \(\mathbb Q[\theta_G,s]\) has \(p+m\leq2m=n(n+1)\) indeterminates.  By \(\mathrm{lem:cubic\mbox{-}graph\mbox{-}ideal\mbox{-}presentation}\), its ideal \(I_G\) is generated in total degree at most three.  The cited Dub\'e bound in \(\mathrm{lem:dube\mbox{-}grobner\mbox{-}degree\mbox{-}bound}\) consequently bounds every element of a reduced Gr\"obner basis, for every lexicographic order, by
\[
 L_n:=2\left(\frac{15}{2}\right)^{2^{n(n+1)-1}}<U_2(n). \tag{2}
\]

It remains to prove that such a basis contains the required linear certificate; this is the extraction step, not part of the cited theorem.  Fix \(q\in\theta_G\), order the other source coordinates ahead of \(q\), and order \(q\) ahead of all covariance coordinates.  Let \(\mathcal B_q\) be the reduced basis for that lexicographic order, and put
\[
 K=I_G\cap\mathbb Q[s],\qquad
 J_q=I_G\cap\mathbb Q[q,s],\qquad
 \mathcal H_q=\mathcal B_q\cap\mathbb Q[q,s],\qquad
 \mathcal H_0=\mathcal B_q\cap\mathbb Q[s].
\]
We first verify the elimination property needed below.  If \(f\in J_q\), the leading monomial of some \(g\in\mathcal B_q\) divides the leading monomial of \(f\).  The latter contains none of the eliminated source variables.  Thus the leading monomial of \(g\) contains none of them.  Under the chosen lexicographic order, any term of \(g\) containing such a variable would be larger than every term containing none, so \(g\in\mathbb Q[q,s]\).  Hence \(\mathcal H_q\) is a Gr\"obner basis of \(J_q\).  The identical argument with \(q\) also eliminated shows that \(\mathcal H_0\) is a Gr\"obner basis of \(K\).

The ring \(A=\mathbb Q[s]/K\) is a domain because it embeds, by covariance substitution, in \(\mathbb Q[\theta_G]\).  Let \(F=\operatorname{Frac}(A)\).  Birationality says that the image of \(q\) belongs to \(F\).  Hence there exist \(a,b\in\mathbb Q[s]\) such that
\[
 a\notin K,\qquad \tau_G(a)q=\tau_G(b),\qquad qa-b\in J_q. \tag{3}
\]
Reduce \(a\) and \(b\) by the covariance-only basis \(\mathcal H_0\).  Subtracting the corresponding elements of \(K\) from \(qa-b\) leaves an element
\[
 f'=q a'-b'\in J_q, \tag{4}
\]
where \(a'\neq0\) is reduced with respect to \(\mathcal H_0\); nonzeroness follows from \(a\notin K\).  Since \(\mathcal H_q\) is a basis of \(J_q\), the polynomial \(f'\) reduces to zero by \(\mathcal H_q\).  No element of \(\mathcal H_0\) can reduce a monomial in the nonzero coefficient \(q a'\), because \(a'\) is already reduced.  Moreover, an element whose degree in \(q\) exceeds one cannot have a leading monomial dividing a monomial of \(q a'\).  Therefore \(\mathcal H_q\setminus\mathcal H_0\) contains an element of the form
\[
 h_q=q a_q(s)-b_q(s). \tag{5}
\]
If \(a_q\in K\), then \(h_q\in I_G\) would also give \(b_q\in K\), and both terms of \(h_q\) would be reducible by \(\mathcal H_0\), contradicting reducedness of \(\mathcal B_q\).  Thus \(a_q\notin K\), equivalently \(\tau_G(a_q)\neq0\).  Equations (2) and (5) give
\[
 \deg h_q\leq L_n<U_2(n). \tag{6}
\]

The argument applies separately to every \(q\in\theta_G\).  Each \(h_q\) uses only \(q\) and \(s\), so these certificates may be placed in any order starting from \(S=\varnothing\).  The definition of \(ID_G\) and (6) yield \(ID_G\leq U_2(n)\).  Maximizing over \(G\in\mathcal G_{2}^{\mathrm{rat}}(n)\) proves \(D_2(n)\leq U_2(n)\).
%%% FIELD k6-statement
The six-vertex graph \(K_6\) is generically rationally identifiable, belongs to \(\mathcal G_{2}^{\mathrm{rat}}(6)\), and satisfies
\[
 ID_{K_6}\leq7.
\]
No matching degree-seven lower bound and no auxiliary-cutset failure claim is asserted.
%%% FIELD k6-proof
Write \(\sigma_{ab}=\Sigma_{ab}\) and
\[
 x_i=\lambda_{i,i+1},\qquad 1\leq i\leq5.
\]
The directed part of \(K_6\) is a chain, so it is acyclic and has maximum indegree one.  For coefficient indices \(1\leq i<j\leq5\) such that the disturbance pair \(\{i+1,j+1\}\) is missing, the corresponding entry of
\(\Omega=(I_6-\Lambda)^{\mathsf T}\Sigma(I_6-\Lambda)\) is
\[
 E_{ij}(x_i,x_j)
 =\sigma_{ij}x_ix_j-\sigma_{i,j+1}x_i-\sigma_{i+1,j}x_j+\sigma_{i+1,j+1}=0. \tag{7}
\]
This identity follows from the covariance parametrization and the support restriction \(\mathrm{ass:omega\mbox{-}support}\); it is a cubic graph-ideal relation.

Choose the shared coefficient
\[
 t=x_3.
\]
For \(k=1,2\), define the covariance-linear polynomials
\[
 D_k(t)=\sigma_{k3}t-\sigma_{k4},\qquad
 N_k(t)=\sigma_{k+1,3}t-\sigma_{k+1,4}.
\]
Then (7) for the missing pairs indexed by \((1,3)\) and \((2,3)\) gives
\[
 D_k(t)x_k-N_k(t)=0,\qquad k=1,2. \tag{8}
\]
For \(k=4,5\), define
\[
 D_k(t)=\sigma_{3k}t-\sigma_{4k},\qquad
 N_k(t)=\sigma_{3,k+1}t-\sigma_{4,k+1}.
\]
The missing pairs indexed by \((3,4)\) and \((3,5)\) similarly give
\[
 D_k(t)x_k-N_k(t)=0,\qquad k=4,5. \tag{9}
\]

Use the remaining missing pair in each triangle to define
\[
 \begin{aligned}
 H(t)&=D_1(t)D_2(t)E_{12}\!\left(\frac{N_1(t)}{D_1(t)},\frac{N_2(t)}{D_2(t)}\right),\\
 J(t)&=D_4(t)D_5(t)E_{45}\!\left(\frac{N_4(t)}{D_4(t)},\frac{N_5(t)}{D_5(t)}\right).
 \end{aligned} \tag{10}
\]
The products in (10) mean the polynomials obtained after clearing the displayed denominators.  Every summand is one covariance coordinate times a product of two covariance-linear polynomials.  Thus
\[
 H(t)=h_2t^2+h_1t+h_0,\qquad
 J(t)=j_2t^2+j_1t+j_0, \tag{11}
\]
where every \(h_j,j_j\in\mathbb Q[s]\) has covariance degree at most three.  Equations (7)--(9), without division, imply \(H(x_3)=J(x_3)=0\) after substitution by \(\tau_{K_6}\).  Consequently
\[
 j_2H(t)-h_2J(t)
 =\Delta(s)t-N(s), \tag{12}
\]
where
\[
 \Delta=j_2h_1-h_2j_1,\qquad
 N=h_2j_0-j_2h_0,\qquad
 \deg_s\Delta,\deg_sN\leq6. \tag{13}
\]
It remains to verify the denominator condition rather than assume it.

Take \(x_i=i/(i+1)\), give \(\Omega\) unit diagonal, give every allowed off-diagonal entry the value \(1/100\), and keep the six graphically missing entries equal to zero.  For every nonzero \(v\in\mathbb R^6\), the inequality \(2|v_iv_j|\leq v_i^2+v_j^2\) gives
\[
 v^{\mathsf T}\Omega v
 \geq \sum_{i=1}^6v_i^2-
 \frac1{100}\sum_{i<j}(v_i^2+v_j^2)
 \geq\frac{95}{100}\sum_{i=1}^6v_i^2>0.
\]
Thus \(\Omega\succ0\), so this is a legal parameter point.  At this point, exact multiplication of \(\Sigma=(I_6-\Lambda)^{-\mathsf T}\Omega(I_6-\Lambda)^{-1}\) gives
\[
 \Sigma=\begin{pmatrix}
 1&51/100&7/20&109/400&57/250&1/5\\
 51/100&63/50&169/200&511/800&263/500&34/75\\
 7/20&169/200&47/30&707/600&361/375&37/45\\
 109/400&511/800&707/600&1509/800&381/250&257/200\\
 57/250&263/500&361/375&381/250&2789/1250&2807/1500\\
 1/5&34/75&37/45&257/200&2807/1500&185/72
 \end{pmatrix}.
\]
Substitution of these entries into (10) gives, up to nonzero rational scalar factors,
\[
 H(t)=(4t-3)(940t-1309),\qquad
 J(t)=(4t-3)(705860t-1145331). \tag{14}
\]
The second roots \(1309/940\) and \(1145331/705860\) are distinct, so the coefficient \(\Delta\) in the leading-term cancellation (12) is nonzero at this legal point.  At the true value \(t=3/4\), the four propagation denominators in (8)--(9) evaluate respectively to
\[
 -\frac1{100},\quad-\frac1{200},\quad-\frac{401}{400},\quad-\frac{401}{500}, \tag{15}
\]
and hence are nonzero as well.  Thus \(\tau_{K_6}(\Delta)\) and each \(\tau_{K_6}(D_k(x_3))\) are nonzero polynomials on the model.  Positivity is used here only to ensure that the nonvanishing witness belongs to the admissible covariance model; a polynomial nonzero at one such point is generically nonzero.

Equations (12)--(13) now give the valid first recursive certificate
\[
 x_3\Delta(s)-N(s)\in I_{K_6},\qquad
 \deg\{x_3\Delta-N\}\leq7, \tag{16}
\]
with a denominator nonzero on the model.  After adjoining \(x_3\), equations (8)--(9) identify \(x_1,x_2,x_4,x_5\).  Each has total certificate degree at most three, since its highest-degree term is \(x_k\sigma_{ab}x_3\), and (15) verifies that its denominator is not the zero polynomial.  Finally, for every free disturbance coordinate,
\[
 \omega_{ab}-\bigl[(I_6-\Lambda)^{\mathsf T}\Sigma(I_6-\Lambda)\bigr]_{ab}=0 \tag{17}
\]
is a degree-at-most-three certificate with denominator one after the structural coefficients have been adjoined.  Hence the displayed order identifies every coordinate of \(\theta_{K_6}\) with maximum degree seven, proving \(ID_{K_6}\leq7\).

The same formulas give a rational inverse to \(\phi_{K_6}\) on the nonempty Zariski-open set where \(\Delta\prod_{k\in\{1,2,4,5\}}D_k(x_3)\neq0\).  Therefore \(\phi_{K_6}\) is generically birational onto its image.  Together with acyclicity and indegree one, this proves \(K_6\in\mathcal G_{2}^{\mathrm{rat}}(6)\).  The proof supplies no lower bound beyond what an exact closure computation may later establish.
%%% FIELD stopping-proof
By \(\mathrm{thm:effective\mbox{-}elimination\mbox{-}upper\mbox{-}bound}\), every \(G\in\mathcal G_{2}^{\mathrm{rat}}(n)\) satisfies
\[
 ID_G\leq U_2(n)=B_2(n). \tag{18}
\]
Moreover \(B_2(n)\geq2\) for every positive integer \(n\).  Apply \(\mathrm{lem:dbi\mbox{-}fixed\mbox{-}cutoff\mbox{-}guarantee}\) with \(d=B_2(n)\).  If \(G\in\mathcal G_{2}^{\mathrm{rat}}(n)\), (18) makes the call return ``yes'', and the adaptive internal search returns formulas at the least degree at which an identifying order succeeds.  Conversely, under the stipulated acyclicity and indegree bound, the same fixed-cutoff guarantee returns ``no'' whenever \(G\) is not rationally identifiable.  Thus ``yes'' is equivalent to membership in \(\mathcal G_{2}^{\mathrm{rat}}(n)\).

At a fixed finite cutoff, the weighted-degree loop contains finitely many monomials, recursive states, and exact row reductions.  Removing the wall-clock timeout therefore leaves a terminating computation.  Substitution of \(d=B_2(n)\) into the fixed-cutoff arithmetic bound gives
\[
 O\!\left(B_2(n)n^5\{3B_2(n)n\}^{4\alpha_{\mathrm{alg}}B_2(n)^2\log(3B_2(n)n)}\right),
\]
as claimed.

For a ``no'' output, retain the final exact rank-intersection transcript.  At every recursive round it records, for each coordinate \(q\) outside the current set \(S\), the row-reduction equality
\[
 \dim A_{B_2(n)}(S)+
 \dim\bigl(qC_{B_2(n)}(S)\bigr)
 =\dim\bigl(A_{B_2(n)}(S)+qC_{B_2(n)}(S)\bigr), \tag{19}
\]
whenever that coordinate is unavailable.  The terminal proper fixed point and (19) are finite rational data.  By \(\mathrm{thm:closure\mbox{-}characterization}\), they certify that no recursive identifying order exists at degree \(B_2(n)\).  In conjunction with (18), such a proper closure is impossible for a member of \(\mathcal G_{2}^{\mathrm{rat}}(n)\); hence the transcript is a complete nonmembership certificate on the stated acyclic indegree-two input domain.  No timeout or unresolved frontier premise enters the argument.  The size of \(B_2(n)\) is only an effective upper bound, and this proof makes no assertion that it is sharp.
%%% FIELD open-kill-partial
The graph \(K_6\) is now proved generically rationally identifiable with \(ID_{K_6}\leq7\): choosing \(x_3=\lambda_{3,4}\) as the shared seed produces two covariance-degree-three quadratics, and their leading-term cancellation gives a degree-seven certificate with a legal positive-definite nonvanishing witness.  Also \(\omega_{11}=s_{11}\), so \(\omega_{11}\in\operatorname{Cl}_1(K_6)\).  No exact lower closure has been certified.  At degree four, after adjoining this root variance, a single availability test compares the span of \(\binom{21+4}{4}=12650\) degree-at-most-four covariance monomial images with a coordinate multiple of the \(\binom{21+3}{3}=2024\) degree-at-most-three images.
%%% FIELD open-frontier-partial
The unchanged full class now has the unconditional effective bound
\[
 3\leq D_2(n)\leq U_2(n)=2\cdot8^{2^{n(n+1)}}\qquad(n\geq3),
\]
with \(D_2(1)=1\), \(D_2(2)=2\), and \(D_2(3)=3\).  The upper bound follows from cubic graph-ideal generators, the cited Dub\'e Gr\"obner-degree bound, and a proved lexicographic extraction of direct linear certificates.  The proposed windmills cannot contribute a lower bound because their source dimension exceeds the observation dimension by \(2\ell-3\), so they are generically infinite-to-one.  Whether \(\sup_nD_2(n)\) is finite, and the sharp constant or growth rate, remains undetermined.
